\documentclass[11pt]{article}

% ---------------------------------------------------------------------------
% Packages
% ---------------------------------------------------------------------------
\usepackage[margin=1in]{geometry}
\usepackage{amsmath,amsthm,amssymb}
\usepackage{graphicx}
\usepackage[numbers,sort&compress]{natbib}
\usepackage{hyperref}
\usepackage{booktabs}
\usepackage{caption}
\usepackage{subcaption}
\usepackage{microtype}
\usepackage{xcolor}

\graphicspath{{../code/figures/}}

% ---------------------------------------------------------------------------
% Theorem environments
% ---------------------------------------------------------------------------
\newtheorem{theorem}{Theorem}[section]
\newtheorem{lemma}[theorem]{Lemma}
\newtheorem{proposition}[theorem]{Proposition}
\newtheorem{corollary}[theorem]{Corollary}
\newtheorem{conjecture}[theorem]{Conjecture}
\theoremstyle{definition}
\newtheorem{definition}[theorem]{Definition}
\newtheorem{example}[theorem]{Example}
\theoremstyle{remark}
\newtheorem{remark}[theorem]{Remark}

% ---------------------------------------------------------------------------
% Macros
% ---------------------------------------------------------------------------
\newcommand{\R}{\mathbb{R}}
\newcommand{\E}{\mathbb{E}}
\newcommand{\Prob}{\mathbb{P}}
\newcommand{\sigmamin}{\sigma_{\min}}
\newcommand{\sigmamax}{\sigma_{\max}}

% ---------------------------------------------------------------------------
% Title
% ---------------------------------------------------------------------------
\title{Notes on the Resolution of the Spielman--Teng Conjecture}
\author{Abhi Wadhwa\\
\textit{University of Southern California}}
\date{Fall 2025}

\begin{document}
\maketitle

% ===================================================================
% ABSTRACT
% ===================================================================
\begin{abstract}
The Spielman--Teng conjecture predicted that for an $n \times n$ random matrix $M$
with iid subgaussian entries, $\Prob(\sigma_n(M) \le \varepsilon n^{-1/2}) \le C\varepsilon + e^{-cn}$.
The best prior results left a gap: the constant $C$ couldn't be pushed to~$1$.
In 2024, Sah, Sahasrabudhe, and Sawhney proved the conjecture up to a $1 + o(1)$ factor,
showing $\Prob(\sigma_n(M) \le \varepsilon n^{-1/2}) \le (1+o(1))\varepsilon + e^{-\Omega(n)}$.
These notes walk through their proof architecture and verify the bound computationally.
\end{abstract}

% ===================================================================
% 1. INTRODUCTION
% ===================================================================
\section{Introduction}
\label{sec:intro}

I saw the Spielman--Teng conjecture for the first time on a half-erased whiteboard in
Zixiang's office hours. He'd been going on about smoothed analysis and the simplex method,
why it works so unreasonably well in practice. In the corner someone had written a question
about the smallest singular value of a random matrix. I didn't think much of it then.
I should have.

For an $n \times n$ matrix $A$, the condition number
\[
  \kappa(A) = \frac{\sigmamax(A)}{\sigmamin(A)}
\]
controls almost everything you'd care about: stability of linear solvers, convergence
of iterative methods, whether your code gives you six good digits or garbage. Worst-case
condition numbers get absurd fast. The $50 \times 50$ Hilbert matrix has
$\kappa(H_{50}) > 10^{17}$. But matrices that come from actual data, or random
perturbations of structured inputs---these behave far better. Why?

Spielman and Teng~\cite{spielman2004smoothed} made this question precise with smoothed
analysis. Start with an adversarially chosen matrix $M$, hit it with a small random
perturbation $\sigma G$, and ask: how large is $\E[\kappa(M + \sigma G)]$? They
conjectured $O(n/\sigma)$, no matter what $M$ you start with. The conjecture sat open
for close to twenty years. Sankar, Spielman, and Teng~\cite{sankar2006smoothed} got
$O(n^{3/2}/\sigma)$. Tao and Vu~\cite{tao2008condition} reached $O(n/\sigma)$ but
with logarithmic losses they couldn't kill.

The sharper question, the one that random matrix theorists actually cared about, was
different. Forget the adversarial base matrix. Let $M$ be random with iid entries.
How likely is it that $\sigma_n(M)$ is tiny?

Edelman~\cite{edelman1988eigenvalues} answered this exactly for Gaussian matrices in
1988. The rescaled quantity $n \cdot \sigma_n(M)$ converges in distribution to a
random variable with density
\begin{equation}
  \label{eq:edelman-density}
  f(x) = (x + 1) e^{-x^2/2 - x}, \quad x \ge 0.
\end{equation}
So $\Prob(\sigma_n(G) \le \varepsilon n^{-1/2}) \approx \varepsilon$ for small
$\varepsilon$, where $G$ is Gaussian. Clean.

For non-Gaussian entries, Rudelson and Vershynin~\cite{rudelson2008littlewood,rudelson2009smallest}
proved
\begin{equation}
  \label{eq:rv-bound}
  \Prob\!\left(\sigma_n(M) \le \frac{\varepsilon}{\sqrt{n}}\right) \le C\varepsilon + e^{-cn},
\end{equation}
where $C > 1$ is an absolute constant. Tao and Vu~\cite{tao2010random} pushed $C$
down to $1$ for $\varepsilon \ge n^{-c}$, but the full range stayed out of reach.
Tikhomirov~\cite{tikhomirov2020singularity} did something else entirely---he pinned
the singularity probability of a Rademacher ($\pm 1$) matrix at exactly
$2^{-n + o(n)}$.

The Spielman--Teng conjecture, in its sharpest form, wanted $C = 1$. Sah,
Sahasrabudhe, and Sawhney~\cite{sah2024spielman} got it, up to a $1 + o(1)$ factor.
That's what I'm writing about here.

The rest goes like this. Section~\ref{sec:prelim} sets up definitions and background.
Section~\ref{sec:main-result} states the main theorems. Section~\ref{sec:proof-arch}
is the core---a walk through the four-step proof.
Section~\ref{sec:experiments} is Monte Carlo experiments checking the bound against
reality. Section~\ref{sec:discussion} is what's still open.

% ===================================================================
% 2. PRELIMINARIES
% ===================================================================
\section{Preliminaries}
\label{sec:prelim}

\subsection{Subgaussian Random Variables}

\begin{definition}
A centered random variable $\xi$ is \emph{subgaussian} with parameter $K$ if
$\E[e^{t\xi}] \le e^{K^2 t^2 / 2}$ for all $t \in \R$. The subgaussian norm is
$\|\xi\|_{\psi_2} = \inf\{K > 0 : \E[e^{\xi^2/K^2}] \le 2\}$.
\end{definition}

Standard Gaussians are subgaussian. So are Rademacher random variables ($\pm 1$
with equal probability), bounded random variables, and anything else whose tails
die off at least as fast as a Gaussian's. Vershynin~\cite{vershynin2018high},
Chapters 2--3, has the full story.

\subsection{Singular Values}

For an $n \times n$ matrix $M$, the singular values $\sigma_1(M) \ge \sigma_2(M) \ge
\cdots \ge \sigma_n(M) \ge 0$ are the eigenvalues of $(M^* M)^{1/2}$. The least
singular value $\sigma_n(M)$ equals the distance from $M$ to the set of
singular matrices:
\[
  \sigma_n(M) = \min_{\|x\| = 1} \|Mx\| = \inf\{t \ge 0 : M - tI \text{ is singular for some rotation}\}.
\]
When $M$ has iid entries of mean~0 and variance~$1/n$, the bulk distribution of
singular values obeys the Mar\v{c}enko--Pastur law~\cite{marchenko1967distribution}.
What happens at the edges---the largest and smallest singular values---is harder
and more interesting.

\subsection{Covering Numbers and $\varepsilon$-Nets}

The $\varepsilon$-net argument is one of those tools you see everywhere in random
matrix theory. An $\varepsilon$-net of the unit sphere $S^{n-1}$ is a finite set
$\mathcal{N}$ where every point on $S^{n-1}$ lies within distance $\varepsilon$
of something in $\mathcal{N}$. The covering number $\mathcal{N}(S^{n-1}, \varepsilon)$
satisfies
\[
  \left(\frac{1}{\varepsilon}\right)^n \le \mathcal{N}(S^{n-1}, \varepsilon)
  \le \left(\frac{2}{\varepsilon} + 1\right)^n.
\]
This turns continuous optimization problems (like bounding $\sigma_n(M)$) into
finite union bounds over a discrete set. Simple idea, absurdly useful. See
Vershynin~\cite{vershynin2018high}, Chapter~4.

% ===================================================================
% 3. THE MAIN RESULT
% ===================================================================
\section{The Main Result}
\label{sec:main-result}

Let $M$ be an $n \times n$ matrix with iid entries drawn from a subgaussian
distribution $\xi$ with mean~$0$ and variance~$1$. Here's what Sah, Sahasrabudhe,
and Sawhney~\cite{sah2024spielman} proved:

\begin{theorem}[Sah--Sahasrabudhe--Sawhney, Theorem 1.1]
\label{thm:sss-main}
For all $\varepsilon \ge 0$,
\begin{equation}
  \label{eq:sss-main}
  \Prob\!\left(\sigma_n(M) \le \frac{\varepsilon}{\sqrt{n}}\right)
  \le (1 + o(1))\varepsilon + e^{-\Omega_\xi(n)},
\end{equation}
where $o(1)$ decays as $C_\xi (\log n)^{-1/16}$.
\end{theorem}

The implied constant in $\Omega_\xi(n)$ and $C_\xi$ depend on the subgaussian
norm of the entry distribution $\xi$---but not on $\varepsilon$, and not on any
detail of the distribution beyond its first two moments and subgaussian parameter.

They also get universality:

\begin{theorem}[Sah--Sahasrabudhe--Sawhney, Theorem 1.2]
\label{thm:sss-universality}
Let $G$ be an $n \times n$ matrix with iid $N(0,1)$ entries. Then
\begin{equation}
  \label{eq:sss-universality}
  \Prob\!\left(\sigma_n(M) \le \frac{\varepsilon}{\sqrt{n}}\right)
  = (1 + o(1))\,\Prob\!\left(\sigma_n(G) \le \frac{\varepsilon}{\sqrt{n}}\right)
  + e^{-\Omega_\xi(n)}.
\end{equation}
\end{theorem}

Since Edelman already worked out the Gaussian case exactly, this gives:

\begin{corollary}[Corollary 1.3]
\label{cor:sss}
For $0 \le \varepsilon \ll 1$,
\begin{equation}
  \label{eq:sss-corollary}
  \Prob\!\left(\sigma_n(M) \le \frac{\varepsilon}{\sqrt{n}}\right)
  = (1 + o(1))\varepsilon + e^{-\Omega_\xi(n)}.
\end{equation}
\end{corollary}

\begin{remark}
The $o(1)$ term decaying as $(\log n)^{-1/16}$ is probably not optimal.
Sah, Sahasrabudhe, and Sawhney note that their argument loses a logarithmic factor
at the truncation step (Section~\ref{sec:truncation}), and improving the exponent $-1/16$
would require tighter control on the regularity events. Whether the $o(1)$ can be
removed entirely---giving $C = 1$ on the nose---remains open.
\end{remark}

% ===================================================================
% 4. PROOF ARCHITECTURE
% ===================================================================
\section{Proof Architecture}
\label{sec:proof-arch}

The proof of Theorem~\ref{thm:sss-main} has four steps. I'll go through each one
and try to get at the ideas without drowning in technical detail. The full argument
lives in Sections 3--9 of~\cite{sah2024spielman}. I spent three weeks reading those
sections. Some of them twice.

\subsection{Step 1: Geometric Reduction}
\label{sec:geometric}

Write $M = \begin{pmatrix} M^* \\ X^T \end{pmatrix}$, where $M^*$ is the
$(n-1) \times n$ matrix consisting of the first $n-1$ rows and $X \in \R^n$
is the last row. The idea: a rank-1 update formula ties $\sigma_n(M)$ to
the geometry of $X$ relative to $M^*$.

\begin{lemma}[Rank-1 update, Lemma 3.2 of~\cite{sah2024spielman}]
\label{lem:rank1}
Let $v \in \ker(M^*)$ be a unit vector (assuming $M^*$ has rank $n-1$, which
happens with high probability). Let $u_1, \ldots, u_{n-1}$ be right singular
vectors of $M^*$ with corresponding singular values $\sigma_1(M^*) \ge \cdots
\ge \sigma_{n-1}(M^*)$. Then
\[
  \sigma_n(M)^2 = \frac{|\langle v, X \rangle|^2}{1 + \sum_{i=1}^{n-1}
  \frac{\langle u_i, X \rangle^2}{\sigma_i(M^*)^2}}.
\]
\end{lemma}

The denominator is $\tilde{\chi}(X)^2$, a weighted sum measuring how much
$X$ lines up with the singular directions of $M^*$, weighted inversely by
singular values. So $\sigma_n(M) \le \varepsilon n^{-1/2}$ becomes
\[
  |\langle v, X \rangle| \le (1 + \varepsilon^{1/4}) \varepsilon n^{-1/2} \tilde{\chi}(X).
\]
Geometrically: the projection of the last row onto the kernel of $M^*$ has to
be small compared to a chi-type statistic.

Why is this useful? It splits the randomness. Condition on $M^*$, and now $v$ and
the singular structure are frozen---all that's left is controlling
$|\langle v, X \rangle|$ relative to $\tilde{\chi}(X)$. That's an
anti-concentration problem, and anti-concentration is something we know how to
think about.

\subsection{Step 2: Truncation}
\label{sec:truncation}

The chi-type statistic $\tilde{\chi}(X)$ sums over all $n-1$ singular directions
of $M^*$. The small singular values blow up the corresponding terms, and those are
a pain to deal with. So you truncate: swap $\tilde{\chi}$ for a truncated version
$\chi$ that only sums over the $\ell = \sqrt{\log n}$ smallest singular directions.

\begin{definition}
Define the truncated chi statistic
\[
  \chi(X)^2 = 1 + \sum_{i=n-\ell}^{n-1} \frac{\langle u_i, X \rangle^2}{\sigma_i(M^*)^2},
\]
where $\ell = \lfloor \sqrt{\log n} \rfloor$.
\end{definition}

Truncation introduces error. To keep it under control, you need a \emph{regularity
event} $R = R_1 \cap R_2 \cap R_3 \cap R_4$ that pins down the singular value
behavior of $M^*$. The four sub-events:
\begin{itemize}
  \item $R_1$: $\sigma_{n-1}(M^*) \ge (\log n)^{-3} n^{-1/2}$. The second-to-last
    singular value isn't too small.
  \item $R_2$: For all $k$, $|\langle v_{n-k}, X \rangle| \le \max\{k^{1/8}, \log\log n\}$
    and $\sigma_{n-k}(M^*) \ge k^{3/4} n^{-1/2}$. The projections and singular values
    are well-behaved at all scales.
  \item $R_3, R_4$: Technical conditions on the $(\log\log n)^2$ scale, keeping
    the truncated chi close to the full chi.
\end{itemize}

Sah, Sahasrabudhe, and Sawhney show (Sections 4--7 of~\cite{sah2024spielman}) that
\begin{equation}
  \label{eq:truncation}
  \Prob\!\left(\sigma_n(M) \le \frac{\varepsilon}{\sqrt{n}} \text{ and } R^c\right)
  = O(\delta_n \varepsilon),
\end{equation}
where $\delta_n = (\log n)^{-c}$ for some constant $c > 0$. The ``bad'' event---regularity
fails and the singular value is small at the same time---costs only $o(\varepsilon)$.

These regularity estimates are the longest part of the paper. They need careful
bounds on singular value gaps of $(n-1) \times n$ random matrices, and the main
input is a bound on $\sigma_{n-k}(M^*)$ giving polynomial-in-$k$ growth away from
the edge. This is where the $(\log n)^{-1/16}$ loss comes from: the exponent
$-1/16$ falls out of balancing truncation level $\ell$, regularity failure
probability, and approximation quality in Lindeberg. I confess I got confused by
the bookkeeping here on my first reading and had to start over.

\subsection{Step 3: Gaussian Replacement}
\label{sec:gaussian-replacement}

This is the heart of the thing. On the regularity event $R$, we've reduced to
\[
  \Prob_X\!\left(|\langle v, X \rangle| \le (1 + \delta_n)\varepsilon n^{-1/2} \chi(X)\right).
\]
The truncated chi $\chi(X)$ depends on $X$ only through $\ell = \sqrt{\log n}$
projections, and regularity keeps those projections under control. What we want: replace
$X$ with a Gaussian vector $Z$ and show the two probabilities are close.

Lindeberg exchange does this. Swap the entries of $X$ for Gaussian entries one at a
time. At each swap, a Taylor expansion bounds the change in probability. Truncation
to $\ell$ directions is what makes the whole thing work---the probability we're
tracking depends on $X$ through $\langle v, X \rangle$ and $\ell$ other projections,
and each replacement shifts these projections by at most $O(1)$.

The specific machinery is ``negative correlation inequalities'' from Campos, Jenssen,
Michelen, and Sahasrabudhe~\cite{campos2023negative}, which handle anti-concentration
of linear forms in weakly correlated random variables. The payoff (Section 8
of~\cite{sah2024spielman}):
\begin{equation}
  \label{eq:lindeberg}
  \Prob_X\!\left(|\langle v, X \rangle| \le (1 + \delta_n)\varepsilon n^{-1/2} \chi(X)\right)
  \le \Prob_Z\!\left(|\langle v, Z \rangle| \le (1 + 2\delta_n)\varepsilon n^{-1/2} \chi(Z)\right)
  + o(\varepsilon).
\end{equation}
The Gaussian side is now tractable. $\langle v, Z \rangle$ is standard Gaussian
(since $v$ is a unit vector and $Z$ has iid $N(0,1)$ entries), and $\chi(Z)$ is
a chi-type random variable independent of $\langle v, Z \rangle$ conditional on
$M^*$.

One thing to flag. The Lindeberg exchange doesn't work on the original (untruncated)
$\tilde{\chi}$ because the full sum hits all $n-1$ singular directions and the
per-step error doesn't decay fast enough. The truncation to
$\ell = \sqrt{\log n}$ directions is exactly what drives the total Lindeberg
error down to $o(\varepsilon)$. I initially thought you could skip the truncation
and just accept a worse rate---you can't, because the error actually diverges.

\subsection{Step 4: Rescaling and Universality}
\label{sec:rescaling}

After replacing with Gaussians, we have a bound in terms of Gaussian singular value
probabilities. The last step takes the truncated-chi formulation back to a singular
value statement.

Set $\varepsilon_0 = n^{-c'}$ for a small constant $c'$. On the Gaussian side,
Tao--Vu universality for row replacement gives
\[
  \Prob\!\left(\sigma_n(\tilde{M}) \le \frac{\varepsilon_0}{\sqrt{n}}\right) + o(\varepsilon),
\]
where $\tilde{M}$ is the matrix $M$ with its last row replaced by a Gaussian row.
Since $\sigma_n(G)$ has the Edelman distribution (which governs the $n^{-1/2}$-scale
behavior), this yields
\[
  \Prob\!\left(\sigma_n(M) \le \frac{\varepsilon}{\sqrt{n}}\right)
  \le (1 + o(1))\,\Prob\!\left(\sigma_n(G) \le \frac{\varepsilon}{\sqrt{n}}\right)
  + e^{-\Omega(n)}.
\]
Plug in $\Prob(\sigma_n(G) \le \varepsilon n^{-1/2}) \le \varepsilon$ from Edelman,
and you're done.

This rescaling step also hands you Theorem~\ref{thm:sss-universality} for free: the
subgaussian probability is $(1 + o(1))$ times the Gaussian probability, so the least
singular value distribution is universal up to the $o(1)$ factor.

\subsection{Summary of the Architecture}

\begin{center}
\fbox{\parbox{0.9\textwidth}{
\textbf{Proof outline:}
\begin{enumerate}
  \item \textbf{Geometric reduction.} Use rank-1 update to reduce $\sigma_n(M)
    \le \varepsilon n^{-1/2}$ to an anti-concentration statement about
    $|\langle v, X \rangle|$ relative to $\tilde{\chi}(X)$.
  \item \textbf{Truncation.} Replace $\tilde{\chi}$ with truncated $\chi$
    (summing over $\sqrt{\log n}$ smallest directions). On a high-probability
    regularity event $R$, the error is $o(\varepsilon)$.
  \item \textbf{Gaussian replacement.} Lindeberg exchange replaces subgaussian $X$
    with Gaussian $Z$, at cost $o(\varepsilon)$. Uses negative correlation
    inequalities.
  \item \textbf{Rescaling.} Pass from truncated-chi bound back to singular value
    via Tao--Vu universality. The Gaussian case gives the optimal constant.
\end{enumerate}
}}
\end{center}

The $1 + o(1)$ loss comes from truncation and Lindeberg, each chipping in a
$(\log n)^{-c}$ error. Getting $C$ exactly to~$1$ would mean either killing the
truncation (hard---the untruncated chi is too wild for Lindeberg) or finding a
proof that doesn't go through ``reduce to anti-concentration, then Gaussianize''
at all. I don't see how to do either, and I suspect neither do they.

% ===================================================================
% 5. COMPUTATIONAL VERIFICATION
% ===================================================================
\section{Computational Verification}
\label{sec:experiments}

I ran Monte Carlo experiments to check the bound against actual numbers. The code
generates random matrices---Gaussian and Rademacher---computes their least singular
values, and stacks the empirical tail probabilities up against
Theorem~\ref{thm:sss-main}. It took about four hours on my laptop, mostly the
$n = 500$ runs.

\subsection{Verifying the Main Bound}

Figure~\ref{fig:sss-bound} plots $\Prob(\sigma_n(M) \le \varepsilon n^{-1/2})$
against $\varepsilon$ for Gaussian matrices at $n = 50, 200, 500$. The bound says
this should stay below $(1 + o(1))\varepsilon$. At all three dimensions, the empirical
CDF hugs the line $y = \varepsilon$. The green dotted line is the
$(1 + \delta_n)\varepsilon$ envelope with $\delta_n = (\log n)^{-1/16}$.
Everything falls below it. Good.

\begin{figure}[htbp]
  \centering
  \includegraphics[width=\textwidth]{sss_bound_verification.png}
  \caption{Empirical CDF of $\sigma_n$ (rescaled) vs the SSS bound
    $(1+o(1))\varepsilon$. Blue: empirical. Red dashed: ideal $\varepsilon$.
    Green dotted: $(1+\delta_n)\varepsilon$ envelope. Based on 2000 trials per dimension.}
  \label{fig:sss-bound}
\end{figure}

\subsection{Universality}

Figure~\ref{fig:universality} puts the CDFs of $\sigma_n$ for Gaussian and Rademacher
matrices side by side. By $n = 200$, the two curves are basically on top of each other.
Theorem~\ref{thm:sss-universality} in action.

\begin{figure}[htbp]
  \centering
  \includegraphics[width=\textwidth]{universality_comparison.png}
  \caption{Universality of the least singular value CDF: Gaussian vs Rademacher.
    The distributions converge rapidly with increasing $n$.}
  \label{fig:universality}
\end{figure}

\subsection{Convergence of the Edelman Density}

Figure~\ref{fig:edelman} shows the empirical distribution of $n \cdot \sigma_n$
against the Edelman limiting density~\eqref{eq:edelman-density}. At $n = 50$ the
fit is already decent. By $n = 500$ it's hard to tell the histogram from the
curve.

\begin{figure}[htbp]
  \centering
  \includegraphics[width=\textwidth]{edelman_comparison.png}
  \caption{Scaled least singular value $n \cdot \sigma_n$ for Gaussian matrices,
    overlaid with the Edelman density $f(x) = (x+1)e^{-x^2/2 - x}$.}
  \label{fig:edelman}
\end{figure}

\subsection{Histograms Across Distributions}

Figure~\ref{fig:histograms} shows raw histograms of $\sigma_n$ for both Gaussian
and Rademacher matrices at six dimensions ($n = 10$ to $500$). As $n$ grows, the
mass piles up near zero and shifts left---consistent with
$\sigma_n = \Theta(1/n)$.

\begin{figure}[htbp]
  \centering
  \includegraphics[width=\textwidth]{least_sv_histograms.png}
  \caption{Histograms of $\sigma_n$ for Gaussian and Rademacher matrices at six
    dimensions. Each histogram is based on 2000 independent trials.}
  \label{fig:histograms}
\end{figure}

\subsection{Convergence Rate of the $o(1)$ Term}

The bound gives $o(1) = O((\log n)^{-1/16})$. Can you see this in
finite-dimensional experiments? Sort of. Figure~\ref{fig:convergence} plots the
ratio $\Prob(\sigma_n \le \varepsilon n^{-1/2}) / \varepsilon$ as a function of
$n$ for several values of $\varepsilon$. This ratio should go to~$1$. It does, but
slowly---which is what you'd expect from a logarithmic rate.

\begin{figure}[htbp]
  \centering
  \includegraphics[width=\textwidth]{convergence_ratio_vs_n.png}
  \caption{Ratio $\Prob / \varepsilon$ vs $n$ for four values of $\varepsilon$.
    Blue: Gaussian. Red: Rademacher. Green dotted: the $1 + (\log n)^{-1/16}$
    upper envelope.}
  \label{fig:convergence}
\end{figure}

Figure~\ref{fig:error-decay} isolates the error term $|\Prob/\varepsilon - 1|$ for
$\varepsilon = 0.3$ and plots it against $n$ on a log-log scale. The theoretical
rate $(\log n)^{-1/16}$ sits above the data as an upper envelope.

\begin{figure}[htbp]
  \centering
  \includegraphics[width=0.75\textwidth]{error_term_decay.png}
  \caption{Decay of the error term $|\Prob/\varepsilon - 1|$ with $n$. Green
    dotted: $(\log n)^{-1/16}$ reference.}
  \label{fig:error-decay}
\end{figure}

\subsection{Tail Probabilities Across Dimensions}

Figure~\ref{fig:prob-bounds} shows $\Prob(\sigma_n < \varepsilon n^{-1/2})$ as
a function of dimension for several $\varepsilon$ values. At every dimension I
tested, the probabilities sit below the Rudelson--Vershynin bound
$C\varepsilon + e^{-cn}$.

\begin{figure}[htbp]
  \centering
  \includegraphics[width=0.75\textwidth]{prob_bounds_vs_n.png}
  \caption{Tail probability vs dimension for Gaussian matrices at several
    $\varepsilon$ levels.}
  \label{fig:prob-bounds}
\end{figure}

\subsection{Condition Number Experiments}

Figures~\ref{fig:cond-scaling}--\ref{fig:cond-theory} check the condition
number side of things. The median condition number of $\sigma G$ scales as
$n/\sigma$ across dimensions (Figure~\ref{fig:cond-scaling}). And the
$50 \times 50$ Hilbert matrix---whose unperturbed condition number is
genuinely terrifying---gets tamed by Gaussian perturbation down to the same
$n/\sigma$ scaling (Figure~\ref{fig:cond-theory}). That fact still strikes me
as a little miraculous.

\begin{figure}[htbp]
  \centering
  \includegraphics[width=0.75\textwidth]{condition_number_scaling.png}
  \caption{Median condition number vs $\sigma$ for Gaussian matrices at several
    dimensions. Dashed gray: $n/\sigma$ reference.}
  \label{fig:cond-scaling}
\end{figure}

\begin{figure}[htbp]
  \centering
  \includegraphics[width=\textwidth]{condition_number_histograms.png}
  \caption{Distribution of $\log_{10}\kappa$ at $\sigma = 0.1$ across dimensions.}
  \label{fig:cond-hist}
\end{figure}

\begin{figure}[htbp]
  \centering
  \includegraphics[width=0.75\textwidth]{condition_vs_theory.png}
  \caption{Median condition number for zero and Hilbert base matrices vs the $n/\sigma$
    bound at $n = 50$.}
  \label{fig:cond-theory}
\end{figure}

% ===================================================================
% 6. DISCUSSION
% ===================================================================
\section{Discussion}
\label{sec:discussion}

The Sah--Sahasrabudhe--Sawhney result closes something that's been open since
Edelman's 1988 density computation, through Rudelson--Vershynin, Tao--Vu,
Tikhomirov. The least singular value of a random matrix with iid subgaussian
entries satisfies $\Prob(\sigma_n \le \varepsilon n^{-1/2}) \le (1 + o(1))\varepsilon$.
It matches the Gaussian prediction up to a vanishing multiplicative factor.

What's left is the $o(1)$ term. It decays as $(\log n)^{-1/16}$. Glacially slow.
Pushing $C$ all the way to $1$ would mean the Rademacher least singular value
distribution matches the Gaussian one with zero error, not just vanishing error.
That's a different kind of statement. The proof as it stands loses the $o(1)$
at truncation and again at Lindeberg exchange. Avoiding both losses would probably
require abandoning the whole ``reduce to anti-concentration, then Gaussianize''
strategy. I don't know what would replace it.

Theorem~\ref{thm:sss-universality}, the universality statement, might actually be
the more surprising result. It says the full distribution of $\sigma_n$ is
asymptotically the same regardless of which subgaussian entries you use. Our
experiments back this up---the Gaussian and Rademacher CDFs are visually identical
by $n = 200$.

And practically? The bound confirms what people doing numerical linear algebra have
known in their bones for decades. Random matrices are well-conditioned. Small
perturbations of structured matrices are well-conditioned. The condition number of
$M + \sigma G$ is $O(n/\sigma)$, and that's tight. Spielman--Teng's smoothed
analysis framework now sits on an essentially optimal foundation. I imagine Spielman
is pleased.

\section*{Acknowledgments}
I'm grateful to Zixiang Zhou for pulling me into this problem and for many
conversations about smoothed analysis and random matrix theory, most of them over
bad coffee.

\bibliographystyle{plainnat}
\bibliography{references}

\end{document}
